\documentclass[a4paper,11pt]{ctexart}

\usepackage[top=1.8cm,bottom=1.8cm,left=2.1cm,right=2.1cm]{geometry}
\usepackage{amsmath,amssymb}
\usepackage{enumitem}
\usepackage{fancyhdr}
\usepackage{lastpage}
\usepackage{arydshln}

\pagestyle{fancy}
\fancyhf{}
\renewcommand{\headrulewidth}{0pt}
\fancyfoot[C]{\small 第~\thepage~页\quad 共~\pageref{LastPage}~页}

\linespread{1.28}
\setlength{\parindent}{0pt}

\newcommand{\tr}{\operatorname{tr}}
\newcommand{\rk}{\mathrm{r}}
\newcommand{\T}{\mathrm{T}}

\begin{document}

\begin{center}
  {\Large\bfseries 2017--2018 学年第二学期期末考试 A 卷参考答案}
\end{center}

\vspace{0.8em}

\textbf{一、判断题（16 分，每题 2 分，共 8 题）}

\begin{enumerate}[label=\arabic*., leftmargin=2em, itemsep=0.65em, topsep=0.5em]
  \item \textbf{【正解】}对

  \textbf{【解析】}若矩阵 $A$ 的所有 $r$ 阶子式都为 $0$，那么 $\rk(A) < r$，因此 $A$ 的所有 $r+1$ 阶子式，如果存在都要为 $0$，否则 $\rk(A) \geqslant r + 1$ 就会出现矛盾.

  \item \textbf{【正解】}错

  \textbf{【解析】}“两个 $n$ 维向量组等价的充分必要条件是它们有相同的秩”这个命题是错误的.

  我们用反例：$\left\{\begin{bmatrix} 1 \\ 0 \\ 0 \\ \vdots \\ 0 \end{bmatrix}_{n\times 1}\right\}$、$\left\{\begin{bmatrix} 0 \\ 1 \\ 0 \\ \vdots \\ 0 \end{bmatrix}_{n\times 1}\right\}$ 是 $n$ 维向量组，且它们的秩都是 $1$，显然它们不等价（不能相互线性表示）.

  \item \textbf{【正解】}错

  \textbf{【解析】}$A = \begin{bmatrix} 0 & 1 \\ 1 & 0 \end{bmatrix}$，$B = \begin{bmatrix} 2 & 0 \\ 0 & 1 \end{bmatrix}$ 是两个初等矩阵，但是 $AB = \begin{bmatrix} 0 & 1 \\ 2 & 0 \end{bmatrix}$ 不是初等矩阵.

  \item \textbf{【正解】}对

  \textbf{【解析】}若 $n$ 阶方阵 $A$ 的行列式 $|A| = 0$，则说明 $\rk(A) < n$，因此 $A$ 的列向量线性相关，所以一定有一列可由其余列线性表示.

  \item \textbf{【正解】}对

  \textbf{【解析】}依题意 $f(A) = a_m A^m + a_{m-1} A^{m-1} + \cdots + a_1 A + a_0 I = 0$，且 $a_m \neq 0, a_0 \neq 0$，那么
  \[
    A(a_m A^{m-1} + a_{m-1} A^{m-2} + \cdots + a_1 I) = -a_0 I.
  \]
  反证法：假设 $A$ 不可逆，那么 $|A| = 0$，所以 $|A|\,|a_m A^{m-1} + a_{m-1} A^{m-2} + \cdots + a_1 I| = |-a_0 I|$，等式左边为 $0$，右边为 $(-a_0)^n$，由于 $a_0 \neq 0$，那么左右两边矛盾，所以 $A$ 可逆.

  \item \textbf{【正解】}错

  \textbf{【解析】}举个反例：假设 $A = [1, 1]$，$\beta = 1$，$r = \rk(A) = 1$，$n = 2$，那么 $AX = \beta$ 的展开形式为 $x_1 + x_2 = 1$，其通解为
  \[
    x = a \begin{bmatrix} 0 \\ 1 \end{bmatrix} + (1 - a) \begin{bmatrix} 1 \\ 0 \end{bmatrix},\quad a \text{ 是任意常数.}
  \]
  所以非齐次线性方程组 $AX = \beta$ 的线性无关的解的个数大于 $n - r = 1$.

  \item \textbf{【正解】}错

  \textbf{【解析】}令 $r = 2, n = 2$，$\alpha_1 = \alpha_2 = \begin{bmatrix} 1 \\ 1 \end{bmatrix}$，显然生成的向量空间 $L(\alpha_1, \alpha_2)$ 的维数为 $1$，而不是 $\min(r, n) = 2$.（思考：题目改为 $r$ 个线性无关的 $n$ 维向量才正确.）

  \item \textbf{【正解】}对

  \textbf{【解析】}设 $n$ 阶方阵 $A, B$ 均为正交矩阵，根据正交矩阵的定义，那么 $AA^{\T} = BB^{\T} = I$，则 $AB(AB)^{\T} = ABB^{\T}A^{\T} = AIA^{\T} = I$，所以 $AB$ 也为正交矩阵.
\end{enumerate}

\vspace{0.8em}

\textbf{二、填空题（20 分，每题 4 分，共 5 题）}

\begin{enumerate}[label=\arabic*., leftmargin=2em, itemsep=0.75em, topsep=0.5em]
  \item \textbf{【正解】}$-3$

  \textbf{【解析】}根据行列式的性质：$(-4, 0, 1, 3)$ 和 $(-2, 5, 1, x)$ 是正交的，因此 $8 + 1 + 3x = 0$，所以 $x = -3$.

  \item \textbf{【正解】}$\begin{bmatrix} 0 & 1 & 1 \\ 1 & 0 & 0 \\ 0 & 0 & 1 \end{bmatrix}$

  \textbf{【解析】}由初等矩阵的性质：$A \begin{bmatrix} 0 & 1 & 0 \\ 1 & 0 & 0 \\ 0 & 0 & 1 \end{bmatrix} = B$，$B \begin{bmatrix} 1 & 0 & 0 \\ 0 & 1 & 1 \\ 0 & 0 & 1 \end{bmatrix} = C$，那么：
  \[
    A \begin{bmatrix} 0 & 1 & 0 \\ 1 & 0 & 0 \\ 0 & 0 & 1 \end{bmatrix} \begin{bmatrix} 1 & 0 & 0 \\ 0 & 1 & 1 \\ 0 & 0 & 1 \end{bmatrix} = C
    \implies Q = \begin{bmatrix} 0 & 1 & 0 \\ 1 & 0 & 0 \\ 0 & 0 & 1 \end{bmatrix} \begin{bmatrix} 1 & 0 & 0 \\ 0 & 1 & 1 \\ 0 & 0 & 1 \end{bmatrix} = \begin{bmatrix} 0 & 1 & 1 \\ 1 & 0 & 0 \\ 0 & 0 & 1 \end{bmatrix}.
  \]

  \item \textbf{【正解】}$1$

  \textbf{【解析】}先求 $A$ 的特征值，由 $\det(A - \lambda I) = 0$，那么可以化简得到：（注：这里的符号 $\det(A - \lambda I)$ 表示方阵 $A - \lambda I$ 的行列式，下同）$[(4 + \lambda)(3 - \lambda) - 8](\lambda - a) = 0$，由于 $\lambda = 1$ 是该式的一个解，而 $\lambda = 1$ 时前面中括号的一项不为 $0$，因此 $(\lambda - a) = 0$，那么 $a = 1$.

  \item \textbf{【正解】}$k \begin{bmatrix} 1 \\ 1 \\ 1 \\ 1 \end{bmatrix} + \begin{bmatrix} 1 \\ -2 \\ -3 \\ 0 \end{bmatrix}$，$k$ 为任意常数

  \textbf{【解析】}根据各行元素和为 $0$，那么 $A \begin{bmatrix} 1 \\ 1 \\ 1 \\ 1 \end{bmatrix} = \mathbf{0}_{5 \times 1}$. 由于 $\rk(A) = 3$，$n = 4$，那么齐次解集的基础解系含 $n - r = 1$ 个线性无关的解向量，因此我们找到了方程的齐次解. $\beta = \alpha_1 - 2\alpha_2 - 3\alpha_3$ 这个条件告诉我们 $[\alpha_1, \alpha_2, \alpha_3, \alpha_4] \begin{bmatrix} 1 \\ -2 \\ -3 \\ 0 \end{bmatrix} = \beta$，所以我们找到了特解. 所以通解为：$k \begin{bmatrix} 1 \\ 1 \\ 1 \\ 1 \end{bmatrix} + \begin{bmatrix} 1 \\ -2 \\ -3 \\ 0 \end{bmatrix}$，$k$ 为任意常数.

  \item \textbf{【正解】}$O$（即 $O_{n \times n}$）

  \textbf{【解析】}对 $A_{n \times n}$ 是一个实对称矩阵，那么 $A$ 一定可相似对角化，所以我们可知：
  \[
    A = P \begin{bmatrix} \lambda_1 & & & \\ & \lambda_2 & & \\ & & \ddots & \\ & & & \lambda_n \end{bmatrix} P^{-1},
    \quad\text{那么 }
    A^m = P \begin{bmatrix} \lambda_1^m & & & \\ & \lambda_2^m & & \\ & & \ddots & \\ & & & \lambda_n^m \end{bmatrix} P^{-1}.
  \]
  由于 $A^m = \mathbf{0}$，那么 $\rk(A^m) = 0$. 根据性质：与可逆矩阵相乘不改变矩阵的秩，那么：
  \[
    \rk\left(\begin{bmatrix} \lambda_1^m & & & \\ & \lambda_2^m & & \\ & & \ddots & \\ & & & \lambda_n^m \end{bmatrix}\right) = \rk(A^m) = 0 \implies \lambda_i = 0\ (i = 1, 2, \cdots, n).
  \]
  又 $A = P \begin{bmatrix} \lambda_1 & & & \\ & \lambda_2 & & \\ & & \ddots & \\ & & & \lambda_n \end{bmatrix} P^{-1}$，所以 $\rk(A) = \rk\left(\begin{bmatrix} \lambda_1 & & & \\ & \lambda_2 & & \\ & & \ddots & \\ & & & \lambda_n \end{bmatrix}\right) = 0$，也就是说 $A = O_{n \times n}$.
\end{enumerate}

\vspace{0.8em}

\textbf{三、（8 分）}

\textbf{【正解】}（利用加边法）
\[
  D_n = \begin{vmatrix}
    a_1 & 1 & 1 & \cdots & 1 \\
    1 & a_2 & 1 & \cdots & 1 \\
    1 & 1 & a_3 & \cdots & 1 \\
    \vdots & \vdots & \vdots & \ddots & \vdots \\
    1 & 1 & 1 & \cdots & a_n
  \end{vmatrix}
  = \widetilde{D}_{n+1}
  = \left|\begin{array}{c:ccccc}
    1 & 1 & 1 & 1 & \cdots & 1 \\ \hdashline
    0 & a_1 & 1 & 1 & \cdots & 1 \\
    0 & 1 & a_2 & 1 & \cdots & 1 \\
    0 & 1 & 1 & a_3 & \cdots & 1 \\
    0 & \vdots & \vdots & \vdots & \ddots & \vdots \\
    0 & 1 & 1 & 1 & \cdots & a_n
  \end{array}\right|.
\]
从第二行开始，每一行都减去第一行，那么：
\begin{align*}
  \widetilde{D}_{n+1} &= \begin{vmatrix}
    1 & 1 & 1 & 1 & \cdots & 1 \\
    -1 & a_1 - 1 & & & & \\
    -1 & & a_2 - 1 & & & \\
    -1 & & & a_3 - 1 & & \\
    \vdots & & & & \ddots & \\
    -1 & & & & & a_n - 1
  \end{vmatrix} \\
  &= \left(1 + \sum_{i=1}^n \frac{1}{a_i - 1}\right) \prod_{i=1}^n (a_i - 1),\quad a_i \neq 1,\ i = 1, 2, \cdots, n.
\end{align*}

\vspace{0.8em}

\textbf{四、（12 分）}

\textbf{【正解】}令 $A = \begin{bmatrix} a + 3 & 1 & 2 \\ a & a - 1 & 1 \\ 3(a + 1) & a & a + 3 \end{bmatrix}$，$[A \mid B] = \left[\begin{array}{ccc:c} a + 3 & 1 & 2 & a \\ a & a - 1 & 1 & a \\ 3(a + 1) & a & a + 3 & 3 \end{array}\right]$（此处 $B = \beta$）.

(1) $\beta$ 能由 $\alpha_1, \alpha_2, \alpha_3$ 唯一的线性表示意味着 $AX = B$ 仅有唯一解. 所以 $\rk(A) = \rk(A \mid B) = 3$，那么 $A$ 可逆，则 $\det(A) \neq 0$. 由于 $\det(A) = a^2(a - 1)$，所以，当 $a \neq 0$ 且 $a \neq 1$ 时，$\beta$ 能由 $\alpha_1, \alpha_2, \alpha_3$ 唯一的线性表示.

(2) 当 $a = 1$ 时，
\[
  A = \begin{bmatrix} 4 & 1 & 2 \\ 1 & 0 & 1 \\ 6 & 1 & 4 \end{bmatrix},\qquad
  [A \mid B] = \left[\begin{array}{ccc:c} 4 & 1 & 2 & 1 \\ 1 & 0 & 1 & 1 \\ 6 & 1 & 4 & 3 \end{array}\right]
  \longrightarrow \left[\begin{array}{ccc:c} 1 & 0 & 1 & 1 \\ 0 & 1 & -2 & -3 \\ 0 & 0 & 0 & 0 \end{array}\right].
\]
此时，$\rk(A) = \rk(A \mid B) < 3$，$\beta$ 能由 $\alpha_1, \alpha_2, \alpha_3$ 线性表示，但表示不唯一. 主元列在第一列和第二列，因此我们将 $x_3$ 看作自由变量，那么
\[
  \begin{cases}
    x_1 = -x_3 + 1, \\
    x_2 = 2x_3 - 3, \\
    x_3 = x_3,
  \end{cases}
  \quad\text{即通解为：}
  X = k \begin{bmatrix} -1 \\ 2 \\ 1 \end{bmatrix} + \begin{bmatrix} 1 \\ -3 \\ 0 \end{bmatrix},\ k \in \mathbb{R}.
\]
此时 $\beta = [\alpha_1, \alpha_2, \alpha_3] X = (1 - k)\alpha_1 + (2k - 3)\alpha_2 + k\alpha_3$（$k \in \mathbb{R}$）.

(3) 当 $a = 0$ 时，
\[
  A = \begin{bmatrix} 3 & 1 & 2 \\ 0 & -1 & 1 \\ 3 & 0 & 3 \end{bmatrix},\qquad
  [A \mid B] = \left[\begin{array}{ccc:c} 3 & 1 & 2 & 0 \\ 0 & -1 & 1 & 0 \\ 3 & 0 & 3 & 3 \end{array}\right]
  \longrightarrow \left[\begin{array}{ccc:c} 1 & 0 & 1 & 0 \\ 0 & 1 & -1 & 0 \\ 0 & 0 & 0 & 1 \end{array}\right].
\]
由于 $\rk(A) \neq \rk(A \mid B)$，因此 $\beta$ 不能由 $\alpha_1, \alpha_2, \alpha_3$ 线性表示.

\vspace{0.8em}

\textbf{五、（10 分）}

\textbf{【正解】}由于 $A = \begin{bmatrix} -2 & 0 & 0 \\ 2 & a & 2 \\ 3 & 1 & b \end{bmatrix}$ 和矩阵 $B = \begin{bmatrix} -1 & & \\ & 2 & \\ & & c \end{bmatrix}$ 相似，那么 $A, B$ 有相同的特征值，则 $\tr(A) = \tr(B)$，$\det(A) = \det(B)$. 由于 $B$ 的特征值为对角线上元素 $-1, 2, c$，那么 $A$ 也有这三个特征值.

由于 $A = \left[\begin{array}{c:cc} -2 & 0 & 0 \\ \hdashline 2 & a & 2 \\ 3 & 1 & b \end{array}\right]$ 是下三角的分块矩阵，所以其特征值为：$-2$，以及 $\begin{bmatrix} a & 2 \\ 1 & b \end{bmatrix}$ 的两个特征值.

那么 $c = -2$.

由于 $\det(B) = -1 \times 2 \times (-2) = 4$，$\tr(B) = -1 + 2 + (-2) = -1$，因此：
\[
  \det(A) = -2 \times (ab - 2) = 4,\qquad \tr(A) = a + b - 2 = -1,
\]
解得：$a = 0, b = 1$（或者 $a = 1, b = 0$）.

\vspace{0.8em}

\textbf{六、（12 分）}

\textbf{【正解】}$f(x, y, z) = x^2 + 4xy + 4xz + y^2 + 4yz + z^2$ 对应的矩阵形式表示为
\[
  f(x, y, z) = [x\ y\ z] \begin{bmatrix} 1 & 2 & 2 \\ 2 & 1 & 2 \\ 2 & 2 & 1 \end{bmatrix} \begin{bmatrix} x \\ y \\ z \end{bmatrix}.
\]
我们令 $A = \begin{bmatrix} 1 & 2 & 2 \\ 2 & 1 & 2 \\ 2 & 2 & 1 \end{bmatrix}$，那么由于 $A$ 是对称矩阵，所以 $A$ 一定可以相似对角化.

令 $\det(A - \lambda I) = 0$，即 $-(\lambda - 5)(\lambda + 1)^2 = 0$，解得 $\lambda_1 = \lambda_2 = -1$，$\lambda_3 = 5$.

当 $\lambda_1 = \lambda_2 = -1$ 时，根据 $(A - \lambda I)x = 0$ 可以解得对应的特征向量为 $p_1 = \begin{bmatrix} 1 \\ -1 \\ 0 \end{bmatrix}$，$p_2 = \begin{bmatrix} 1 \\ 0 \\ -1 \end{bmatrix}$. 将其进行施密特正交化，那么
\[
  q_1 = p_1 = \begin{bmatrix} 1 \\ -1 \\ 0 \end{bmatrix},\qquad
  q_2 = p_2 - \frac{q_1 \cdot p_2}{q_1 \cdot q_1} q_1 = \begin{bmatrix} 1 \\ 0 \\ -1 \end{bmatrix} - \frac{1}{2}\begin{bmatrix} 1 \\ -1 \\ 0 \end{bmatrix} = \begin{bmatrix} 1/2 \\ 1/2 \\ -1 \end{bmatrix}.
\]
单位化后为：$\varepsilon_1 = \begin{bmatrix} 1/\sqrt{2} \\ -1/\sqrt{2} \\ 0 \end{bmatrix}$，$\varepsilon_2 = \begin{bmatrix} 1/\sqrt{6} \\ 1/\sqrt{6} \\ -2/\sqrt{6} \end{bmatrix}$.

同理：$\lambda_3 = 5$ 对应的标准化后的特征向量为 $\varepsilon_3 = \begin{bmatrix} 1/\sqrt{3} \\ 1/\sqrt{3} \\ 1/\sqrt{3} \end{bmatrix}$.

那么所使用的正交变换矩阵 $P = \begin{bmatrix} 1/\sqrt{2} & 1/\sqrt{6} & 1/\sqrt{3} \\ -1/\sqrt{2} & 1/\sqrt{6} & 1/\sqrt{3} \\ 0 & -2/\sqrt{6} & 1/\sqrt{3} \end{bmatrix}$（$P$ 正交）.

因此 $A = P \begin{bmatrix} -1 & & \\ & -1 & \\ & & 5 \end{bmatrix} P^{-1} \implies P^{\T} A P = \begin{bmatrix} -1 & & \\ & -1 & \\ & & 5 \end{bmatrix}$（利用 $P^{-1} = P^{\T}$）.

我们令：$X = \begin{bmatrix} x \\ y \\ z \end{bmatrix}$，$V = \begin{bmatrix} v_1 \\ v_2 \\ v_3 \end{bmatrix}$，$X = PV$，所以
\[
  f(x, y, z) = X^{\T} \begin{bmatrix} 1 & 2 & 2 \\ 2 & 1 & 2 \\ 2 & 2 & 1 \end{bmatrix} X
  = V^{\T} P^{\T} \begin{bmatrix} 1 & 2 & 2 \\ 2 & 1 & 2 \\ 2 & 2 & 1 \end{bmatrix} PV
  = V^{\T} \begin{bmatrix} -1 & & \\ & -1 & \\ & & 5 \end{bmatrix} V.
\]
因此，标准型为 $f = -v_1^2 - v_2^2 + 5v_3^2$.

\vspace{0.8em}

\textbf{七、（10 分）}

\textbf{【证明】}根据 $A, B, A + B$ 都是 $n$ 阶可逆矩阵，因此
\[
  A^{-1}(A + B)B^{-1} = A^{-1}AB^{-1} + A^{-1}BB^{-1} = B^{-1} + A^{-1}.
\]
因此 $A^{-1} + B^{-1}$ 可逆，且
\[
  (A^{-1} + B^{-1})^{-1} = \{A^{-1}(A + B)B^{-1}\}^{-1} = B(A + B)^{-1}A.
\]

\newpage

\textbf{八、（10 分）}

\textbf{【证明】}由于矩阵 $A$ 的列向量组 $\{\alpha_1, \alpha_2, \cdots, \alpha_n\}$ 线性无关，因此 $m \geqslant n$. 否则当 $m < n$ 时，$\rk(A) \leqslant \min(m, n) \implies \rk(A) < n$，那么 $A$ 的列向量组 $\{\alpha_1, \alpha_2, \cdots, \alpha_n\}$ 线性相关，故矛盾. 由于 $m \geqslant n$，且 $A$ 的列向量组 $\{\alpha_1, \alpha_2, \cdots, \alpha_n\}$ 线性无关，因此 $\rk(A) = n$.

\textbf{充分性：}(1) $K$ 是可逆矩阵时，那么由于 $B = AK$，因此 $\rk(B) = \rk(A) = n$，因此矩阵 $B$ 的列向量组 $\{\beta_1, \beta_2, \cdots, \beta_n\}$ 也线性无关（注：与可逆矩阵相乘不改变矩阵的秩）.

\textbf{必要性：}(2) 当矩阵 $B$ 的列向量组 $\{\beta_1, \beta_2, \cdots, \beta_n\}$ 线性无关时，那么 $\rk(B) = n$. 因此，我们可以得到 $\rk(B) = \rk(A) = n$. 由题意存在 $B = AK$，注意到 $K$ 是 $n$ 阶方阵.

我们用反证法：假设 $K$ 不可逆，那么一定有 $\rk(K) < n$，因此
\[
  \rk(AK) \leqslant \min(\rk(A), \rk(K)) = \rk(K) < n
\]
（注：秩的不等式 $\rk(AB) \leqslant \min(\rk(A), \rk(B))$）.

又因为 $\rk(B) = n$，$B = AK$，所以 $\rk(AK) = n$，因此产生了矛盾，也就是说 $K$ 一定是可逆矩阵.

\end{document}
